% TOP_PROBLEM: {"id": 2302079, "title": "Research Problems in Function Theory — Problem 2.79", "category_id": 8, "category": {"id": 8, "name": "analysis", "display_name": "Analysis", "description": "Limits, continuity, calculus, and function theory.", "slug": "analysis", "order_index": 8, "created_at": "2026-07-31T15:26:25.671Z"}, "status": "partially_solved"}
% FIRST_POSTED: null
% ENTRY_KIND: statement_only
% Imported from: batch12.tex; UnsolvedMath ID 2302079
% Compile twice: lualatex 2302079.tex
\documentclass[11pt]{article}
\usepackage{fontspec}
\setmainfont{Latin Modern Roman}
\usepackage{amsmath,amssymb,mathtools,unicode-math}
\setmathfont{Latin Modern Math}
\usepackage{xurl,hyperref}
\hypersetup{colorlinks=true,urlcolor=blue,pdftitle={UnsolvedMath Problem 2302079}}
\newcommand{\sref}[2]{\hyperlink{source:#1:#2}{[S#2]}}
\newcommand{\cataloguescope}[1]{\paragraph{Mathematical question.}#1}
\begin{document}
% UnsolvedMath ID 2302079; source catalogue code AMR-022-2079
\setcounter{section}{2302078}
\section{Research Problems in Function Theory — Problem 2.79}\label{2302079}
{\small\sffamily UnsolvedMath ID 2302079 · Analysis}
\subsection{Definitions and mathematical statement}

\paragraph{Shared notation.}$\mathbb N=\{1,2,\ldots\}$, and $\mathbb Z,\mathbb Q,\mathbb R,\mathbb C$ denote the integers, rationals, reals and complex numbers. Any other conventions are specified locally.


Let $g:\widehat{\mathbb C}\to\widehat{\mathbb C}$ be a rational map of degree at least $2$. Its Julia set $J(g)$ is the complement of the normality set of its iterates. Area means spherical two-dimensional area. A Beltrami differential on $J(g)$ is a measurable coefficient $\mu$ with $\|\mu\|_\infty<1$. It is invariant if, in local complex coordinates,
\[\mu(z)=\mu(g(z))\frac{\overline{g'(z)}}{g'(z)}\]
for almost every $z\in J(g)$; the finitely many critical points do not affect this condition.
The source update reports Buff--Ch\'eritat quadratic counterexamples to part (b), with positive-area Julia sets.
\paragraph{Statement recorded in the supplied source.}
Assume $J(g)\ne\widehat{\mathbb C}$. (a) Must every invariant Beltrami differential on $J(g)$ vanish almost everywhere? (b) Must $J(g)$ have area zero? \sref{2302079}{2}
\subsection{Short English statement}
For a rational map whose Julia set is not the whole sphere, must every invariant Beltrami differential vanish, and must the Julia set have zero area?
\subsection{Sources}
\begin{description}
\item[\hypertarget{source:2302079:1}{S1}] ULAM.AI and the UnsolvedMath Contributors, UnsolvedMath: A Curated Collection of Open Mathematics Problems, record 2302079 (AMR-022-2079). CC BY 4.0. \url{https://huggingface.co/datasets/ulamai/UnsolvedMath}.
\item[\hypertarget{source:2302079:2}{S2}] W. K. Hayman and E. F. Lingham, \emph{Research Problems in Function Theory}, arXiv:1809.07200v2 (2018), Problem 2.79 and Update 2.79. \url{https://arxiv.org/abs/1809.07200}.
\end{description}
\label{end:2302079}
\end{document}
